\lysh{34159}{2}

\begin{alist}
\item Η συνάρτηση ορίζεται για $x \in \mathbb{R}$ με:
\[x - 3 \ne 0 \Leftrightarrow x \ne 3.\]
Άρα το πεδίο ορισμού της $f$ είναι το $A = \mathbb{R} - \{3\}$.

\item Θα παραγοντοποιήσουμε το $x^2 - 5x + 6$.
Το τριώνυμο $x^2 - 5x + 6$ έχει $\alpha = 1, \beta = -5, \gamma = 6$ και διακρίνουσα:
\[\Delta = \beta^2 - 4\alpha\gamma = (-5)^2 - 4 \cdot 1 \cdot 6 = 25 - 24 = 1 > 0.\]
Οι ρίζες του τριωνύμου είναι οι:
\[x_{1,2} = \dfrac{-\beta \pm \sqrt{\Delta}}{2\alpha} = \dfrac{-(-5) \pm \sqrt{1}}{2 \cdot 1} = \dfrac{5 \pm 1}{2} = \begin{cases} 3 \\ 2 \end{cases}\]
Τότε:
$x^2 - 5x + 6 = (x - 2)(x - 3)$.
Για $x \ne 3$ ο τύπος της $f$ γράφεται:
\[f(x) = \dfrac{(x - 2)(x - 3)}{x - 3} = x - 2.\]

\item Για τις τετμημένες των σημείων τομής της $C_f$ με τον άξονα $x'x$ λύνουμε την εξίσωση:
\[f(x) = 0 \Leftrightarrow x - 2 = 0 \Leftrightarrow x = 2.\]
Άρα η $C_f$ τέμνει τον άξονα $x'x$ στο σημείο $A(2, 0)$.
Επίσης έχουμε:
$f(0) = 0 - 2 = -2$.
Άρα η $C_f$ τέμνει τον άξονα $y'y$ στο σημείο $B(0, -2)$.
\end{alist}
